%september 8
Notation:
$\phi_{a}(\rho) := \rho \circ(\nvf(F,\rho) \to a)$, 
$\phi_{a_1,\dots, a_t} = \phi_{a_2,\dots,a_t}(\phi_{a_1}(\rho))$.

Fix $a \in \cube{t}$.
Let $R_1 = \bra{\rho:F|_{\phi_{a_1,\dots,a_t}(\rho)} = 0, F|_{\phi_{a_1,\dots,\bar{a_t}}(\rho)} = 1 }$,
$R_2 = \bra{\rho: F|_{\phi_{a_2,\dots,a_t}(\rho)} = 0, F|_{\phi_{a_2,\dots,\bar{a_t}}(\rho)} = 1}$.

Consider the communication matrix $M$ whose rows and columns are indexed by $X$ and $Y$ respectively. Where $X = \bra{\phi_{a_1}(\rho): \rho \in R_1}$ and $Y = \bra{\rho : \phi_{\bra{a_1}(\rho)}: \rho \in R_1}$. The protocol $\prot$ partitions $M$ into combinatorial rectangles which correspond to transcripts.  Let $A_{\ts_1} \times B_{\ts_1}$ be the rectangle which corresponds to the transcript $\ts_1$. 

\begin{claim}
    \label{calim: prob rectangles}
        Let $M$ be the communication matrix whose rows and columns are indexed by $X$ and $Y$ respectively. For each transcript $\ts_1$, let
        $\pr{A_{\ts_1}} = \sum_{\rho \in A_{\ts_1}}$, $\pr{B_{\ts_1}} = \sum_{\rho \in B_{\ts_1}}$.
%        $X = \cup_{\ts_1 \in \prot} A_{\ts_1}$, $Y = \cup_{\ts_1 \in \prot} B_{\ts_1}$.
        \[
        \sum_{\ts_1 \in \prot} \pr{A_{\ts_1}} \pr{B_{\ts_1}} \leq \pr{X} \pr{Y}
        \]
                
\end{claim}

\begin{claim}
    \[
    \pr{\rho: \rho \in R_1} \leq \sqrt{\prot} \sqrt{\pr{X} \pr{Y}} \overset{?}{\leq} \sqrt{|\prot|}\pr{\rho: \rho \in R_2}
    \]
\end{claim}

\begin{proof}
\begin{align*}
    %\pr{\cdt(F, \rho) has depth =t}
    \pr{\rho: \rho \in R_1} %&= \sum_{a} \pr{\rho: \rho \in R_1(\bar{a})}\\
                                    &=  \sum_{x_{i_1}}\pr{\rho: \nvf(F,\rho) = i_1, \rho \in R_1}\\
                                    &\leq \sum_{x_{i_1}} \sum_{\ts_1}\pr{\rho: \prot(\phi_{a_1}(\rho), \phi_{\bar{a_1}}(\rho)) = \ts_1, \rho \in R_1}\\
    \text{Let } E :=\bra{\rho \in R_1: \prot(\phi_{a_1}(\rho), \phi_{\bar{a_1}}(\rho)) = \ts_1} \\
                                    &\leq \sum_{x_{i_1}} \sum_{\ts_1} \sum_{\rho \in E} \pr{\rho}\\
                                    &\leq O(p) \sum_{x_{i_1}} \sum_{\ell_1}\sum_{\rho \in E} \sqrt{\pr{\phi_{a_1}(\rho)}}\sqrt{\pr{\phi_{\bar{a_1}}(\rho)}}\\
                                    &\leq O(p) \sum_{x_{i_1}, \ts_1} \sqrt{\sum_{\rho \in E}\pr{\phi_{a_1}(\rho)}} \sqrt{\sum_{\rho \in E}\pr{\phi_{\bar{a_1}}(\rho)}}\\
                                    &= O(p) \sum_{x_{i_1},\ts_1} \sqrt{\pr{A_{\ts_1}} \pr{B_{\ts_1}}}\\
                                    &\leq O(p) \sqrt{\sum_{x_{i_1},\ts_1} 1 } \sqrt{\sum_{x_{i_1},\ts_1}\pr{A_{\ts_1}} \pr{B_{\ts_1}}}\\
                                    &= \sqrt{|\prot|} \sqrt{\pr{X} \pr{Y}}
\end{align*}

\end{proof}

\subsection{Decision-tree size after random restriction}

Let $R(s) = \bra{\rho: \cdt(F,\rho) \text{ has size } = s}$,\\ 
$R(s_1,s_2) = \bra{\rho: \cdt(F,\phi_{a_1}(\rho)) \text{ has size } s_1 \text{ and } \cdt(F,\phi_{bar{a}}(\rho)) \text{ has size } = s_2}$.\\
$X(s_1,s_2) = \bra{\phi_{a_1}(\rho) : \rho \in R(s_1,s_2)}$,\\
$Y(s_1,s_2) = \bra{\phi_{\bar{a_1}}(\rho): \rho \in R(s_1,s_2)}$.\\
\begin{claim}
    $$\prob{\rho}{\cdt(F,\rho) \text{has size = } s} \leq (s \cdot p\sqrt{\prot})^{\log s}$$
\end{claim}
\begin{proof}
\begin{align*}
   \pr{R} &= \sum_{s_1 = 1}^{s-1}
                \sum_{x_{i_1}} 
                    \sum_{\ts_1} \pr{\rho: \nvf(F,\rho) = i_1, \rho \in R(s_1,s_2)}\\
        &= \sum_{s_1 = 1}^{s-1}
            \sum_{x_{i_1}} 
                \sum_{\ts_1} \pr{\rho: \prot( \phi_{a_1}(\rho),                \phi_{\bar{a_1}}(\rho) ) = \ts_1, \rho \in R(s_1,s_2)}\\
    \text{Let } E :=\bra{\rho \in R(s_1,s_2): \prot(\phi_{a_1}(\rho), \phi_{\bar{a_1}}) = \ts_1}, \\                
        &= \sum_{s_1 = 1}^{s-1}
            \sum_{x_{i_1}}        
                \sum_{\ts_1} \sum_{\rho \in E} \pr{\rho}\\
        &= O(p) \sum_{s_1 = 1}^{s-1}
                \sum_{x_{i_1}}
                \sum_{\ts_1} \sum_{\rho \in E} \sqrt{ \pr{\phi_{a}(\rho)}} \sqrt{\pr{\phi_{\bar{a}}(\rho)} } \\
        &= O(p) \sum_{s_1 = 1}^{s-1}
                \sum_{x_{i_1}}
                \sum_{\ts_1} \sqrt{\sum_{\rho \in E} \pr{\phi_{a}(\rho)}} \sqrt{\sum_{\rho \in E} \pr{\phi_{\bar{a}(\rho)}}}\\
        &= O(p) \sum_{s_1 = 1}^{s-1}
            \sum_{x_{i_1}} \sum_{\ts_1} \sqrt{\pr{A_{\ts_1}} \pr{B_{\ts_1}}}\\
        &\leq O(p) \sum_{s_1 = 1}^{s-1}
                \sqrt{\sum_{x_{i_1}} \sum_{\ts_1} 1 } \sqrt{\sum_{x_{i_1}} \sum_{\ts_1} \pr{A_{\ts_1}} \pr{B_{\ts_1}}} \\
        & = O(p \sqrt{|\prot|}) \sum_{s_1 =1}^{s-1}\sqrt{\pr{X(s_1,s_2)} \pr{Y(s_1,s_2)}} \\
        &= O(p \sqrt{|\prot|}) \sum_{s_1 = 1}^{s -1}\sqrt{\pr{R(s_1)} \pr{R(s -s_1)}}\\
        &= O\cbra{\cbra{p \sqrt{2s|\prot|}}^{\log s}} \tag*{by Claim \ref{claim: recursion}}
\end{align*}
\end{proof}

\begin{claim}
    \label{claim: recursion}
    $$f(s) = \alpha \sum_{s_1 = 1}^{s-1} \sqrt{f(s_1) f(s-s_1)}$$
    $$\Rightarrow f(s) \leq (\alpha \sqrt{2s})^{\log s}$$
\end{claim}


\begin{proof}
    Suppose $f(s_1) \leq \cbra{\sqrt{2s_1} \cdot \alpha}^{\log s_1}$ for all $s_1 < s$. 

    $f(s_1) f(s-s_1) = \cbra{\sqrt{2s_1} \alpha}^{\log s_1} \cbra{\sqrt{(s -s_1)} \alpha}^{\log (s-s_1)}$\\
    
    \begin{align*}
        \log \cbra{f(s_1) f(s-s_1)} 
        &= \log s_1(\log \sqrt{2s_1}\alpha) + \log (s-s_1) \cbra{\log \sqrt{2(s-s_1)}\alpha}\\
        &= \frac{\log^2 s_1 + \log^2 (s-s_1)}{2} + \cbra{\log \sqrt{2}\alpha}                                                   
        \cbra{\log s_1 + \log (s-s_1)}
    \end{align*}
Consider the one variable real function, 
    $g(x) = \log x (\log \sqrt{2}\alpha) + \frac{\log^2 {x}}{2}$ where $\sqrt{2x}\alpha \leq 1$.\\
    
    $g'(x) = \frac{\log \sqrt{2}\alpha}{x} + \frac{\log x}{x}$.
    
    
\end{proof}



\begin{proof}
    Let $s$ be even and suppose $f(s_1)(f(s-s_1))$ achieves its maximum at $s_1 = s/2$.We get the following recursion, 
    \begin{align*}
        f(s) &\leq \alpha \cdot s \cdot f(s/2)\\
            &\leq \alpha^{\log s} 2^{(\log s)(1+ \log s)/2}\\
            &= \alpha^{\log s} \cbra{2^{(1+ \log s)/2}}^{\log s}\\
            &= (\alpha \sqrt{2s})^{\log s}. 
    \end{align*}
\end{proof}